\section{Conclusion}

The raw material for primitive machine life may not be a humanoid shell or a
single extraordinary model. It may be the existing, inelegant ecology of
servers, operating systems, compilers, build chains, caches, repositories,
gates, networks, credentials, executable extensions, and intermittent
intelligent agents.

Those materials become one candidate subject only when organization binds
them. KFD supplies a possible semantic kernel: facts do not drift, direction
outlives actions, authority remains bounded, work and history retain causal
identity, and successor state requires admission. Recursive dogfood supplies
a possible self-production loop: the current organization can develop,
qualify, admit, and then use changes to its own organs. Agent cognition and KFX
execution make that loop intelligent and effective, while Buildchain and
ordinary machines give it a body and metabolism.

\factlabel\ Current Kungfu engineering contains several native mechanisms
needed by this architecture.

\inferencelabel\ Their composition is sufficient to justify building a bounded
proto-organism experiment now.

\hypothesislabel\ A subject organized this way can survive agent, organ,
toolchain, and machine replacement; maintain itself for months; and recursively
improve bounded capabilities with decreasing human repair.

\nonclaimlabel\ No current evidence establishes consciousness, intrinsic
desire, biological life, unrestricted autonomy, or a completed self-sustaining
Kungfu organism.

The appropriate next step is therefore neither proclamation nor dismissal. It
is construction: create a bounded habitat, establish one founding Initiative,
run the heartbeat loop, retain exact evidence, inject failures, measure human
intervention, and allow the result to falsify the claim. If the system survives
while preserving identity and governing its own material change, machine life
will have become less a metaphor and more an engineering phenomenon.
